\chapter{Continuous Hamiltonian cochains}
\label{ch:vol2-continuous-cochains}

A formal Hamiltonian is determined by its finite jets. An operation on
the quantized algebra need not return as many coefficients as it inspects.
For example, extraction of a cubic coefficient can have a scalar
output. A limit of self-valued cochains at equal input and output
jet orders would miss this operation.

We therefore keep the two orders separate. The resulting continuous
Hochschild complex has braces and a relative bar resolution. Its
Chevalley comparison is compatible with both jet orders and with
the quantization parameter. This constructs a continuous algebraic
centre. Relating it to geometric collision operations or to bulk
observables requires maps preserving those additional operations.

\section{Residue duality and opposite directions of variance}

Let $k$ be a field of characteristic zero. The polynomial formulas of
Chapter~\ref{ch:vol2-hamiltonian-jets} apply over $k$. Put
\[
 \mathfrak m=(z_1,z_2),\qquad
 \h=\mathfrak m^2\subset k\llbracket z_1,z_2\rrbracket,
 \qquad \h_{\leq d}=\h/\mathfrak m^{d+1}\quad(d\geq2).
\]
The bracket is \eqref{eq:vol2-hamiltonian-bracket}. Give $\h$ its
order topology and $k$ the discrete topology. A continuous linear
functional on $\h$ vanishes on $\mathfrak m^{d+1}$ for some $d$.
Consequently
\begin{equation}\label{eq:vol2-continuous-dual}
 \h^\vee_{\mathrm{cont}}
   =\varinjlim_d\h_{\leq d}^\vee
   =\bigoplus_{a+b\geq2}k\rho_{ab},
 \qquad
 \rho_{ab}=z_1^{-a-1}z_2^{-b-1}\,dz_1\,dz_2.
\end{equation}
The notation on the right represents the iterated residue functional,
normalized by
\[
 \operatorname{Res}_0(z_1^iz_2^j\rho_{ab})=\delta_{ia}\delta_{jb}.
\]
On the whole power-series ring, the same formula includes all
$a,b\geq0$. The two degree-one modes pair with the translations in
$\h^{\mathrm{aff}}$, whereas the pointed algebra starts in degree two.
The continuous dual has the discrete topology.

\begin{proposition}[Coadjoint currents]\label{prop:vol2-continuous-coadjoint}
For $f=z_1^iz_2^j\in\h$,
\begin{equation}\label{eq:vol2-residue-coadjoint}
 \operatorname{ad}_f^*\rho_{ab}
  =\bigl(j(a+1)-i(b+1)\bigr)\rho_{a-i+1,b-j+1}.
\end{equation}
A mode with a negative index or with total degree less than two is
zero as a functional on $\h$. The formula extends to formal $f$,
and the coadjoint action is continuous.
\end{proposition}

\begin{proof}
By definition,
$(\operatorname{ad}_f^*\rho_{ab})(g)=-\rho_{ab}(\{f,g\})$.
For $g=z_1^pz_2^q$, the monomial bracket has coefficient $iq-jp$
and exponent $(i+p-1,j+q-1)$. The residue can be nonzero only at
$(p,q)=(a-i+1,b-j+1)$, where its negative coefficient is
$j(a+1)-i(b+1)$. This proves the formula.
If $a+b\leq d$, a Hamiltonian of order greater than $d$ acts by
zero: the possible output mode would have degree at most one.
Thus the action on any finite sum of residue modes uses only
finitely many Taylor coefficients of $f$. Its output is again a
finite sum. This also proves continuity for the stated topologies.
\end{proof}

The graded semidirect Lie algebra
\[
 \mathfrak q_{\mathrm{cont}}
   =\h\ltimes\h^\vee_{\mathrm{cont}}[-1]
\]
has degree-zero Hamiltonians and degree-one covectors. The second
summand is abelian. The order topology on the first summand and the
discrete topology on the second make it separated and complete.
Evaluation gives the invariant pairing between the two summands:
its invariance is the defining negative-transpose identity above.

The quotient $\h\to\h_{\leq d}$ induces an inclusion of duals.
Its image is a coadjoint submodule, and there is a surjective Lie map
\begin{equation}\label{eq:vol2-current-subquotient}
 \h\ltimes\h_{\leq d}^\vee[-1]
   \longrightarrow
 \mathfrak q(\h_{\leq d})
   =\h_{\leq d}\ltimes\h_{\leq d}^\vee[-1],
 \qquad (x,\alpha)\longmapsto(x\bmod\mathfrak m^{d+1},\alpha).
\end{equation}
Its kernel is $\mathfrak m^{d+1}$ in the first summand. Indeed,
that ideal acts trivially on the displayed finite dual, and the
coadjoint formula descends to the finite quotient. Thus the finite
current algebra is a quotient of a subalgebra of
$\mathfrak q_{\mathrm{cont}}$.

Simultaneous coordinate truncation on both full summands is not a
Lie map. At $d=2$, it kills both $z_1^3$ and $\rho_{21}$, but
\[
 \operatorname{ad}_{z_1^3}^*\rho_{21}=-6\rho_{02}
\]
has a nonzero quadratic component. The quotient of Hamiltonians and
the inclusion of covectors must retain their opposite directions.

\section{The completed algebra and its continuous operations}

Set $R=k\llbracket\hbarq\rrbracket$ and
$R_N=R/(\hbarq^{N+1})$. For $d\geq2$ and $N\geq0$, define
\begin{equation}\label{eq:vol2-continuous-algebra}
 A_{d,N}=U_{R_N}(\h_{\leq d}\otimes_kR_N,\hbarq\{-,-\}),
 \qquad A=\widehat U_\hbar(\h)=\varprojlim_{d,N}A_{d,N}.
\end{equation}
The defining relation is $xy-yx=\hbarq\{x,y\}$. Choose one ordering
of the Hamiltonian monomials, compatible with the inclusions of
their finite bases. The PBW theorem proved in
Proposition~\ref{prop:vol2-rees-pbw} gives free $R_N$-modules with
the corresponding ordered-monomial bases. Every quotient map in
\eqref{eq:vol2-continuous-algebra} is onto.

Give $A$ the inverse-limit topology with discrete quotients $A_{d,N}$.
At each fixed $(d,N)$, PBW expressions are finite polynomials.
In particular, this construction does not complete the polynomial
degree in the enveloping generators. For $n\geq1$, put
\begin{equation}\label{eq:vol2-quotient-tensor}
 A^{\widehat\otimes_R n}
  =\varprojlim_{d,N}A_{d,N}^{\otimes_{R_N}n},
 \qquad A^{\widehat\otimes_R0}=R.
\end{equation}
Different cutoffs in the tensor factors give the same limit because
a common larger cutoff is cofinal. All continuity statements below
refer to this quotient tensor.

Here is an operation requiring different jet orders. Let $X_{30}$
be the enveloping generator represented by the Hamiltonian $z_1^3$.
Projection to $A_{3,N}$ followed by extraction of the PBW coefficient
of the single monomial $X_{30}$ is an $R_N$-linear scalar-valued map.
The compatible maps in $N$, followed by the scalar inclusion into
$A$, define a continuous $R$-linear map $\lambda_{30}:A\to A$.
Its output in $A_{2,N}$ is nonzero on $X_{30}$, although the
input projection to $A_{2,N}$ kills $X_{30}$. That output therefore
cannot be computed from the quadratic input jet.

\begin{definition}\label{def:vol2-uniform-cochains}
The continuous Hochschild cochains for the quotient tensor are
\begin{equation}\label{eq:vol2-two-index-cochains}
 C^n_{\mathrm{uc}}(A/R,A)
   =\varprojlim_{d,N}\varinjlim_{e\geq d}
       \operatorname{Hom}_{R_N}
       (A_{e,N}^{\otimes_{R_N}n},A_{d,N}).
\end{equation}
The direct maps precompose with input quotients. The inverse maps
reduce the output and the parameter precision. Give each direct
limit the discrete topology and each $C^n_{\mathrm{uc}}$ its
inverse-limit topology. The subscript specifies uniform continuity
for \eqref{eq:vol2-quotient-tensor}.
\end{definition}

\begin{lemma}[An input cutoff for each output]\label{lem:vol2-continuous-maps}
The quotient projections in \eqref{eq:vol2-quotient-tensor} are onto,
and there is a natural identification
\begin{equation}\label{eq:vol2-continuous-hom}
 C^n_{\mathrm{uc}}(A/R,A)
  =\operatorname{Hom}_{R,\mathrm{cont}}
       (A^{\widehat\otimes_R n},A).
\end{equation}
Both the tensor source and the cochain space are separated and complete.
\end{lemma}

\begin{proof}
A finite PBW tensor lifts by the same ordered monomials to all
larger cutoffs; hence every projection is onto. On underlying modules,
symmetrization identifies the tensor source with
\[
 \left(\varprojlim_d
      (\operatorname{Sym}_k\h_{\leq d})^{\otimes_k n}\right)
       \llbracket\hbarq\rrbracket.
\]
Each parameter coefficient is a compatible system of finite
polynomials. Coefficientwise division by $\hbarq^{N+1}$ therefore
identifies its quotient by that parameter power with
$\varprojlim_d A_{d,N}^{\otimes n}$. The inverse-limit description
also gives separatedness and completeness.

Let $f$ be a continuous map on the right of
\eqref{eq:vol2-continuous-hom}. Its composite with the output
$(d,N)$ kills the kernel of some input projection $(e,M)$.
Increase that projection so that $e\geq d$ and $M\geq N$.
Surjectivity makes the induced map $R$-linear. Since the target is
killed by $\hbarq^{N+1}$, it factors through
$A_{e,N}^{\otimes n}$. Two choices give the same map precisely
after a common input refinement. The output compatibilities give
\eqref{eq:vol2-two-index-cochains}.
Conversely, its compatible families define a map into $A$, and their
finite input factorizations at each output make it continuous.
Finally, an inverse limit of discrete modules is a closed submodule
of their product and is complete, proving the assertion for cochains.
\end{proof}

\section{Composition and the relative derived centre}

All enveloping generators have cohomological degree zero, and an
$n$-cochain has cohomological degree $n$. Its Hochschild differential is
\begin{align*}
 (bf)(a_1,\ldots,a_{n+1})={}&a_1f(a_2,\ldots,a_{n+1})\\
 &+\sum_{i=1}^n(-1)^i
       f(a_1,\ldots,a_i a_{i+1},\ldots,a_{n+1})\\
 &+(-1)^{n+1}f(a_1,\ldots,a_n)a_{n+1}.
\end{align*}
For brace signs, use the reduced degree $|f|=n-1$.
If $g_j$ has arity $q_j$, define
\begin{equation}\label{eq:vol2-continuous-braces}
 f\{g_1,\ldots,g_r\}
   =\sum_{\text{ordered disjoint insertions}}
       (-1)^{\sum_j(q_j-1)i_j}
       f(\ldots,g_1(\ldots),\ldots,g_r(\ldots),\ldots),
\end{equation}
where $i_j$ is the number of arguments preceding the block of
$g_j$ in the resulting argument list. Equivalently, these are the
Koszul insertion signs on the suspended tensor coalgebra.
For $f$ of arity $p$ and $g$ of arity $q$, the cup product is
\[
 (f\smile g)(a_1,\ldots,a_{p+q})
   =f(a_1,\ldots,a_p)g(a_{p+1},\ldots,a_{p+q}).
\]

\begin{theorem}[Continuous braces]\label{thm:vol2-continuous-braces}
The differential, cup product and every finite brace preserve
$C^\bullet_{\mathrm{uc}}(A/R,A)$. They are continuous for the
stated cochain topologies and obey the Hochschild brace identities.
\end{theorem}

\begin{proof}
Fix an output quotient $(d,N)$ for the outer cochain $f$. It uses
some input cutoff $e$. Each inner cochain $g_j$ is then needed only
at output $(e,N)$, and that output uses some finite input cutoff
$e_j$. Choose $E$ at least as large as $e,e_1,\ldots,e_r$.
Every summand is now a map from a tensor power of $A_{E,N}$ to
$A_{d,N}$. There are finitely many summands at fixed arities, and
further refinement leaves their class unchanged. This constructs
the brace. Multiplication is defined at each finite quotient, so
the same argument constructs $b$ and the cup product.

To prove continuity at a specified tuple of cochains, fix its outer
component at $(d,N)$ and the finitely many inner components just
used. These equalities describe open neighborhoods, and all outputs
at $(d,N)$ agree throughout them. The argument applies to every
output quotient.
For a brace identity, choose one common quotient for all the
intermediate outputs in its two sides. Both sides enumerate the
same nested or disjoint insertions there. The Koszul sign is the
sign of the same permutation of suspended inputs, so the terms
agree. Thus the identity holds at every output and in the limit.
Associativity of multiplication gives $b^2=0$.
\end{proof}

An individual rectangular complex with coefficients in $A_{d,N}$
and inputs in $A_{e,N}$ is not a self-valued brace algebra when
$e\ne d$. The total cochain space supplies the intermediate output
needed by the next insertion.

We now specify the derived category for this construction. Work with
complete separated linearly topologized $R$-modules and quotient
completed tensors. A map out of a tensor is continuous when each
discrete output factors through discrete quotients of all its inputs.
Use complete $A$-bimodules whose action is continuous for this tensor,
and write $A^e=A\widehat\otimes_R A^{\mathrm{op}}$.
A short exact sequence is \emph{relatively admissible} when it has
a continuous $R$-linear splitting, which need not be $A$-bilinear.
Relative resolutions and derived Hom refer to this exact structure.

For a complete coefficient module $E$ and a complete $A$-bimodule $M$,
restriction to
$1\otimes E\otimes1$ gives the induction adjunction
\begin{equation}\label{eq:vol2-continuous-induction}
 \operatorname{Hom}_{A^e,\mathrm{cont}}
       (A\widehat\otimes_R E\widehat\otimes_R A,M)
   =\operatorname{Hom}_{R,\mathrm{cont}}(E,M).
\end{equation}
The inverse sends $f$ to $a\otimes e\otimes b\mapsto af(e)b$.
First refining the output for the action and then the input for $f$
proves its continuity. If $M'\to M$ has a continuous $R$-linear
section, a map on the right lifts by composition with that section.
The adjunction then lifts the bimodule map. Thus induced modules
are relatively projective.

\begin{theorem}[The completed relative bar]\label{thm:vol2-continuous-bar}
The augmented complex
\[
 B_n(A)=A\widehat\otimes_R A^{\widehat\otimes_R n}
                   \widehat\otimes_R A\quad(n\geq0),
 \qquad B_0(A)\longrightarrow A,
\]
with alternating multiplication boundary and $B_n$ in degree $-n$,
is a relative projective resolution. Hence
\begin{equation}\label{eq:vol2-continuous-centre}
 \operatorname{RHom}^{\mathrm{rel,cont}}_{A^e}(A,A)
    \simeq C^\bullet_{\mathrm{uc}}(A/R,A).
\end{equation}
\end{theorem}

\begin{proof}
Each term is induced, hence relatively projective by
\eqref{eq:vol2-continuous-induction}. On the augmented bar, put
\[
 s_n(a_0\otimes\cdots\otimes a_{n+1})
   =1\otimes a_0\otimes\cdots\otimes a_{n+1},
 \qquad s_{-1}(a)=1\otimes a.
\]
These maps are continuous and $R$-linear. In $\partial s+s\partial$,
every multiplication except the newly inserted first one occurs
twice with opposite signs. The remaining term is the identity.
On cycles, $s$ is therefore a continuous $R$-linear section of the
boundary onto those cycles. This proves admissible exactness.
Continuous bimodule Hom identifies the term in degree $n$ with
\eqref{eq:vol2-continuous-hom}. Its induced differential is exactly
$b$, proving \eqref{eq:vol2-continuous-centre}.
\end{proof}

This is the derived centre in the stated relative continuous theory.
An absolute derived category of topological modules or an analytic
tensor product requires a comparison with this resolution.

\section{The Chevalley comparison and the inverse limit}

For $e\geq d$, let $A_{d,N}$ be an $A_{e,N}$-bimodule through the
quotient. Define
\begin{equation}\label{eq:vol2-rectangular-chevalley}
 P^n_{e,d,N}
  =\operatorname{Hom}_k\bigl(\Lambda^n\h_{\leq e},
                  \operatorname{Sym}_k\h_{\leq d}\bigr)\otimes_kR_N.
\end{equation}
The coefficient action of $\h_{\leq e}$ factors through
$\h_{\leq d}$ and extends its adjoint action as a derivation of
the symmetric algebra. Let $d_{\mathrm{CE}}$ be the Chevalley
differential with this action and with the Poisson Lie bracket.
For example,
\[
 (d_{\mathrm{CE}}f)(x,y)=x\cdot f(y)-y\cdot f(x)-f(\{x,y\})
\]
for a one-cochain. In general the action term indexed by $i$ has
sign $(-1)^i$ on inputs $x_0,\ldots,x_n$, and the bracket term
indexed by $i<j$ has sign $(-1)^{i+j}$, with $\{x_i,x_j\}$ put
first and the remaining inputs in their original order.
The differential on $P^\bullet_{e,d,N}$ is
$\hbarq d_{\mathrm{CE}}$.

\begin{theorem}[Chevalley comparison with two jet orders]
\label{thm:vol2-continuous-chevalley}
Output PBW symmetrization and input antisymmetrization define
natural quasi-isomorphisms
\begin{align}
 C^\bullet_{\HH}(A_{e,N}/R_N,A_{d,N})
   &\longrightarrow (P^\bullet_{e,d,N},\hbarq d_{\mathrm{CE}}),
       \label{eq:vol2-rectangular-map}\\
 f&\longmapsto\left[(x_1,\ldots,x_n)\longmapsto
     \operatorname{PBW}^{-1}\sum_{\sigma\in S_n}
       \operatorname{sgn}(\sigma)
       f(x_{\sigma(1)},\ldots,x_{\sigma(n)})\right].\notag
\end{align}
They commute with input refinement, output quotient and parameter
reduction. Consequently they induce a quasi-isomorphism
\begin{equation}\label{eq:vol2-continuous-chevalley-map}
 C^\bullet_{\mathrm{uc}}(A/R,A)
   \longrightarrow
 \mathcal P_{\mathrm{uc}}
   :=\varprojlim_{d,N}\varinjlim_{e\geq d}
                   (P^\bullet_{e,d,N},\hbarq d_{\mathrm{CE}}).
\end{equation}
For both systems in this formula, the ordinary inverse limit
computes the derived inverse limit.
\end{theorem}

\begin{proof}
Use the finite Chevalley bimodule resolution from the proof of
Theorem~\ref{thm:vol2-relative-rees}, now for $A_{e,N}$ and with
coefficient bimodule $A_{d,N}$. Its terms are
\[
 A_{e,N}\otimes_{R_N}\Lambda^n(\h_{\leq e}\otimes_kR_N)
          \otimes_{R_N}A_{e,N}.
\]
They are free bimodules. Filtering by PBW degree plus exterior
degree gives the Koszul resolution of the diagonal of a polynomial
algebra over $R_N$. The diagonal differences form a regular sequence
even over this nilpotent coefficient ring. A cycle has finite
filtration degree, and lifting an associated-graded primitive lowers
that degree. This proves exactness without a convergence assumption.

Antisymmetrization of the exterior factors is a map into the bar
resolution: internal bar products combine in pairs to give the
scaled Lie commutator, and the two endpoint terms give the left and
right actions. It is therefore a chain homotopy equivalence between
projective resolutions. Applying bimodule Hom into $A_{d,N}$ gives
\eqref{eq:vol2-rectangular-map}. Under PBW, the commutator action of
$x$ on the coefficient algebra is $\hbarq$ times its adjoint Poisson
action; the bracket term also carries $\hbarq$. This proves the
stated differential. All formulas commute with the three quotient
maps, proving naturality.

Filtered direct limits of modules are exact. For the output inverse
limit, use the countable cofinal sequence $(d,N)=(r+2,r)$.
Every cochain at one output factors through a finite input order
$e$. For the next output, increase $e$ if necessary and precompose
with the input quotient. Lift the value of each PBW tensor basis
element to the next output and parameter precision, then extend
linearly over the new coefficient ring. Each value is a finite
polynomial, so this defines an allowed cochain. The exterior basis
gives the same construction on $P^n_{e,d,N}$. These are degreewise
lifts, not chain sections; no $R_N$-linear section of
$R_{N+1}\to R_N$ is used.

Thus both cochain towers are degreewise surjective. For a countable
tower of complexes with transitions $p_r:K_{r+1}\to K_r$, the
derived inverse limit is the homotopy fibre of
\[
 1-\mathrm{shift}:\prod_rK_r\longrightarrow\prod_rK_r.
\]
Degreewise surjectivity of the transition maps makes this product
map degreewise onto: choose the next coordinate recursively to
solve $x_r-p_rx_{r+1}=y_r$. Its kernel is $\varprojlim_rK_r$.
The homotopy fibre is consequently quasi-isomorphic to that kernel.
Applying the derived inverse limit to the levelwise comparisons
therefore gives \eqref{eq:vol2-continuous-chevalley-map}.
\end{proof}

At $\hbarq=0$, the Chevalley differential vanishes, as required by
the commutative polynomial special fibre. At positive parameter
precision the two Chevalley terms remain scaled together.
The theorem identifies complexes. It does not identify a chosen
polyvector wedge with the Hochschild cup product or make
antisymmetrization a strict brace map.

Finite comparison homotopies can also be computed. At parameter
zero, write $t_i=x_i-y_i$ for the diagonal differences in the
Koszul resolution and $\epsilon_i$ for its exterior generators.
On positive total $t,\epsilon$ degree, the operators
\[
 d_0=\sum_i t_i\partial_{\epsilon_i},\qquad
 h_0=\frac{\sum_i\epsilon_i\partial_{t_i}}
                    {\deg_t+\deg_\epsilon}
\]
satisfy $d_0h_0+h_0d_0=1$; on $t,\epsilon$-constant terms the
homotopy is zero. Let $i_0$ include this constant summand and let
$p_0$ project to it. The anticommutator before division is the
total Euler operator, so
\[
 d_0h_0+h_0d_0=1-i_0p_0,\qquad
 h_0^2=h_0i_0=p_0h_0=0.
\]
After PBW identification the quantum differential is
$d_K=d_0+\delta$, where $\delta$ has order $\hbarq$.
Modulo $\hbarq^{N+1}$, put
\[
 T=\sum_{j=0}^{N}(-\delta h_0)^j,\qquad
 h_K=h_0T,\qquad p_K=p_0T,\qquad
 i_K=\sum_{j=0}^{N}(-h_0\delta)^j i_0.
\]
These sums are finite. Multiplying them and using
$d_K^2=0$ and the preceding identities gives
\[
 d_Kh_K+h_Kd_K=1-i_Kp_K,\qquad p_Ki_K=1.
\]
The differential transferred to the constant summand is zero by
degree. That summand represents the diagonal quotient $A_{e,N}$,
so $p_K$ is its augmentation in PBW coordinates. This supplies
the contraction on the augmented resolution.

Starting with the augmentation lift, define the bar-to-Chevalley
comparison on a free middle generator by
\[
 \psi_n(1[a_1|\cdots|a_n]1)
   =h_K\psi_{n-1}\partial(1[a_1|\cdots|a_n]1),
\]
and extend bimodule-linearly. Its argument is a cycle by the previous
stage, so the contraction identity proves the chain-map equation.
If $\alpha$ is the antisymmetrization map into the bar, define on
its free middle generators
\[
 H_n=s_n\bigl((1-\alpha\psi)_n-H_{n-1}\partial\bigr),
 \qquad H_{-1}=0,
\]
and again extend bimodule-linearly. The expression in parentheses
is a cycle by induction, so $\partial H+H\partial=1-\alpha\psi$.
These are finite algorithms at every fixed cutoff.
Transporting higher operations through a simultaneous continuous
choice requires the compatibility of those homotopies as well.

\section{A localization test for limits and adjoints}

The surjectivity used above is a substantive hypothesis. A scalar
example shows both what can fail without it and why changing the
completion changes the problem.

Let $S=k[t]$, $L=k[t,t^{-1}]$, and write $D(S)$ and $D(L)$ for
the derived categories of modules. Restriction of scalars
$i:D(L)\to D(S)$ is fully faithful, since localization is flat and
$L\otimes_S L=L$. Its left adjoint is
$L\otimes_S^{\mathbb L}(-)$; its right adjoint is
$\operatorname{RHom}_S(L,-)$. The right adjoint represents maps
from localized objects into a given object, whereas the left
adjoint represents maps out of that object into localized ones.

\begin{proposition}[The two localization adjoints]
\label{prop:vol2-localization-adjoints}
Their values at $S$ are
\[
 L\otimes_S^{\mathbb L}S=L,
 \qquad
 \operatorname{RHom}_S(L,S)
     \simeq(k\llbracket t\rrbracket/k[t])[-1].
\]
The latter equivalence is one of $S$-complexes. Multiplication by
$t$ is invertible on its displayed cohomology and gives the
$L$-module structure. For $\widehat S=k\llbracket t\rrbracket$
and $\widehat L=k((t))$, the right adjoint takes $\widehat S$ to
zero, while the left adjoint takes it to $\widehat L$.
\end{proposition}

\begin{proof}
The direct system $S\xrightarrow{t}S\xrightarrow{t}\cdots$ has
colimit $L$, where the $n$th copy maps by multiplication by $t^{-n}$.
Its free telescope resolution has terms $\bigoplus_{n\geq0}S$
and differential $e_n\mapsto e_n-te_{n+1}$. Injectivity follows
by inspecting the least nonzero coordinate of a finitely supported
vector. Its cokernel imposes exactly the relations of that colimit.
Applying $\operatorname{Hom}_S(-,S)$ gives the complex in degrees
zero and one
\begin{equation}\label{eq:vol2-polynomial-telescope}
 \prod_{n\geq0}S\xrightarrow{D}\prod_{n\geq0}S,
 \qquad D(a)_n=a_n-ta_{n+1}.
\end{equation}
If $D(a)=0$, each $a_n$ is divisible by every power of $t$, so
$a_n=0$. Define
\[
 \Theta(b)=\left[\sum_{n\geq0}t^nb_n\right]
               \in k\llbracket t\rrbracket/k[t].
\]
The series converges coefficientwise. For $b=D(a)$ it telescopes
to the polynomial $a_0$, so $\Theta$ annihilates the image of $D$.
It is onto, by taking the $b_n$ to be scalar coefficients of a series.
If its value is zero, the series equals a polynomial $c$. Put
\[
 a_n=t^{-n}\left(c-\sum_{j<n}t^jb_j\right).
\]
The expression in parentheses is a polynomial divisible by $t^n$,
so $a_n\in S$, and $D(a)=b$. This proves the asserted cokernel.
On that cokernel, the inverse of multiplication by $t$ sends
$[f]$ to $[(f-f(0))/t]$.

Over $\widehat S$, the same telescope differential has the inverse
\[
 a_n=\sum_{j\geq0}t^j b_{n+j}.
\]
Each series converges $t$-adically, and the homogeneous equation
still has only the zero solution. Thus the derived right-adjoint
value is zero. Flat localization computes both left-adjoint values.
\end{proof}

For $b_n=1$, the class $\Theta(b)=[(1-t)^{-1}]$ is nonzero.
Every finite subsystem nevertheless has a polynomial solution:
with $a_{M+1}=0$, take
\[
 a_n=1+t+\cdots+t^{M-n}\qquad(0\leq n\leq M).
\]
These degrees grow with $M$ and give no polynomial solution of the
infinite system. Completion supplies such a solution by admitting
a different class of coefficients. The continuous Hamiltonian
comparison avoids this inference by proving surjectivity of its
actual cochain transitions.

\section{Compatible models and geometric actions}

The algebra $A$ of \eqref{eq:vol2-continuous-algebra} is the formal
Hamiltonian boundary algebra, and \eqref{eq:vol2-continuous-centre}
constructs its relative continuous derived centre. The explicit
model $\mathcal P_{\mathrm{uc}}$ records the corresponding
Hamiltonian cochains with their two variances. A different
polyvector realization must be compared with this complex in its
own topology.

To compare another realization obtained from finite models, one
must also specify its transitions. At a fixed jet order let
\[
 A_d=\varprojlim_NA_{d,N},\qquad
 \mathcal P_d=
 \bigl(C^\bullet(\h_{\leq d},\operatorname{Sym}\h_{\leq d})
              \llbracket\hbarq\rrbracket,\hbarq d_{\mathrm{CE}}\bigr).
\]
A \emph{compatible jet comparison} consists of complete $R$-linear
cochain models $P_d,H_d$, continuous transition maps, levelwise comparisons
$P_d\simeq\mathcal P_d$ and
$H_d\simeq C^\bullet_{\HH,\mathrm{cont}}(A_d/R,A_d)$, and
compatible quasi-isomorphisms $P_d\to H_d$. All comparison maps
and limit identifications are continuous. Require degreewise
surjective transitions or use derived inverse limits. Specify the
identifications of their chosen limits with a proposed complex
$\mathcal P_\infty$ and with $C^\bullet_{\mathrm{uc}}(A/R,A)$.
An ordinary limit may be used only when its comparison with the
derived limit is supplied. All identifications use the quotient
topology fixed above.

\begin{theorem}[Continuous comparison from compatible jets]
\label{thm:vol2-compatible-jets}
A compatible jet comparison induces a continuous quasi-isomorphism
\[
 \Phi_\infty:\mathcal P_\infty\xrightarrow{\sim}
                C^\bullet_{\mathrm{uc}}(A/R,A).
\]
Compatible homotopy-$\Etwo$ and Swiss-cheese refinements, whose
operations converge in the chosen topology, give the corresponding
complete open--closed operations.
\end{theorem}

\begin{proof}
The derived inverse limit preserves levelwise equivalences. For
degreewise surjective towers it agrees with the ordinary limit by
the product calculation in the proof of
Theorem~\ref{thm:vol2-continuous-chevalley}. Apply the specified
limit identifications. For the final assertion, the supplied
operadic comparisons identify each operation and its coherence
homotopies; convergence defines them on the limits.
\end{proof}

The two-index construction already gives the algebraic cochain
comparison \eqref{eq:vol2-continuous-chevalley-map}. It supplies
neither a chosen geometric polyvector product nor its equivalence
with the brace operations. In particular, an invariant continuous
Duflo map requires a multiplicative comparison and a specified
prescription for its wheel coefficients.

Finite matrix branes are representation-theoretic images of the
formal Hamiltonian boundary algebra whenever their representations
are continuous for this completion. Its relative derived centre is
then an algebraic target for coherent bulk actions. A field-theory
identification starts instead from that theory's observables and
constructs their action map. It must preserve the insertion supports,
collisions, vacuum and higher substitutions. The continuous algebraic
resolution provides the target of this comparison; it does not
construct its geometric or physical source.
